The paper that you hold in your hands (or perhaps all you're holding is a
retinal pattern) has been a long time in coming. On the surface, it is the
revision of a 1996 manuscript containing a long and detailed proof. This
itself is the incarnation of work that goes back some 15 years earlier. In
fact, since as far back as I can remember, Kay Pechenick (as she was then
known, before she became Kay Pechenick DeVicci, before she became Kay
Pechenick DeVicci Shultz) has been fascinated by the problem of determining
the largest cube that would fit inside a tesseract. The question was first
posed by Martin Gardner in his November 1966 {\it Scientific American} column,
but the roots of the problem go back centuries. My tongue was only lightly
planted in cheek when I wrote the following talk abstract for a 2013
conference in honor of Gregory Galperin on his 60th birthday:

\begin{addmargin}[0.4in]{0.0in}
{\bf The De Vicci Cube, and other mysteries from the fourth dimension
(and beyond)}

There is a secret cube that lives in four dimensions called the De Vicci
Cube. ``Surely you mean a `hypercube' or a `tesseract'?'' you might be
forgiven to ask. No, I mean a regular three-dimensional cube. But this is no
ordinary cube. Leonardo Da Vinci didn't know of its existence, nor did the
Freemasons, and nor does Dan Brown (for, if he did, there would surely be
another bestseller and movie out). It was first foreshadowed by the
Babylonians long ago. Then, in the seventeenth century, a certain Prince
Rupert of the Rhine wagered on the outcome an optimization problem. He
got his answer, and he won his bet, but he didn't solve the problem. The
problem wasn't solved until Pieter Nieuwland examined it over 100 years
later. The solution was

1.06066017177982128660126654315...

which, {\it \`{a} la mani\`{e}re \'{e}gyptienne},

= 1 + 1/17 + 1/545 + 1/561804 + ...

Nieuwland was getting close to the Cube, but soon after his discovery, he
died under very mysterious circumstances. And still the De Vicci Cube
remained a secret, lurking in the sea of undiscovered things. What is the De
Vicci Cube? We know that De Vicci's Cube has linear dimension

1.00743475688427937609825359524...

= 1 + 1/135 + 1/36564 + …

This was calculated by at least six people (all dead now) before K. De Vicci
discovered it. The Cube has acquired De Vicci's name because she
supplied an actual proof. Erdös knew about the Cube, but soon after
learning of it, he died too. There is much more to say. More than this, you
must attend the talk, as I have written too much already.
\end{addmargin}

The history reads like a screenplay: It's got Babylonians, ne'er-do-well
princes and kings, untimely death, Latin manuscripts, and the 4th dimension
-- all in one package. Add to this the mystery of a little-known mathematical
genius who lives in Cherry Hill, New Jersey... what more could one want? But
unlike a Hollywood movie, the DeVicci Cube story is actually true.

Some historical notes and corrections are in order. The phrase ``DeVicci
Cube'' derives from Steven R. Finch's coinage ``the DeVicci Constant''
(= 1.0074347...) in his book Mathematical Constants. Also, a study of Martin
Gardner's extensive correspondence reveals that only three others had
calculated the DeVicci Constant before DeVicci. They are Eugen Bosch (1966),
G. de Josselin de Jong (1971), and Hermann Baer (1974). In addition, although
his solution for the 3-cube in a 4-cube is incorrect, Andrew L. Clarke
(in 1967) had already formulated the correct inequalities and other results
for the m-cube in an n-cube generalized Rupert problem.

I first got involved with the problem when, after first hearing from Kay
about her work, I was able to calculate the DeVicci Constant to six decimal
places.  Later I found a proof for the statement f(m, m+1) > 1 (in the
notation of the present paper), and more recently, Terry Ligocki and I have
used techniques from computational physics to deeply probe the full f(m,n)
problem. It is true that both Paul Erd\"{o}s (and Raphael Robinson) looked at
the f(m,n) problem, but neither of them made any inroads into the problem.

In one of his earlier collections of columns, Mathematical Carnival, Gardner
claimed that he had received seven solutions to his four-dimensional Rupert
problem, but that it was difficult for him to evaluate them as none of them
agreed with each other. Gardner re-reviewed the situation in 2001, in his The
Colossal Book of Mathematics, where he devotes a page (page 172) to his
generalization of the Rupert problem. He corrects himself there, but then
goes on to mention two people, Kay R. Pechenick DeVicci and Kay R. Pechenick,
not realizing that these are one and the same person.

Finally, many people are familiar with the generalized Prince Rupert problem
from the one-page description in the book by Croft, Falconer and Guy Unsolved
Problems in Geometry, which is both fortunate in terms of publicity, but
unfortunate in terms of scholarship. Their entry on this problem was
selectively paraphrased from personal letters sent to Richard Guy, they
credit neither Pechenick DeVicci Shultz nor myself, and they get important
details wrong -- a flawed reference.

The present manuscript is also not perfect: first and foremost, its length is
daunting. But it is a primary reference, and it should be viewed as a source
for future work on the Rupert problem.
